\section{Related Work}
\vspace{-0.1in}
\noindent\textbf{Large Language Model Compression.}
The field of model compression for Large Language Models (LLMs) has seen a surge of innovative techniques aimed at mitigating the substantial computation and memory requirements these models demand~\citep{zhu2023survey,yuan2024llmunveiled}. 
Various methods have emerged to address this challenge, each taking a unique approach to reduce the memory footprint of LLMs. These methods primarily fall into three categories: weight quantization~\citep{courbariaux2015binaryconnect,dettmers2022llm}, network pruning~\citep{lecun1989optimal,frantar2023sparsegpt}, and knowledge distillation~\citep{hinton2015distilling,agarwal2023gkd}. 
For the wide body of research on LLM compression, please refer to
\citep{zhu2023survey} for the comprehensive survey. Among these methods, weight quantization has gained significant traction in the context of LLMs due to its effectiveness. However, despite its popularity as a neural network compression technique, low-rank factorization has not been extensively explored in the realm of LLMs. 
Recognizing this gap, we introduce a novel low-rank decomposition method tailored specifically for decomposing the weight matrices of LLMs in a training-free manner.
% This approach aims to provide an efficient and effective solution for LLM compression, leveraging the untapped potential of low-rank factorization in this domain.

\noindent\textbf{Low-rank Decomposition.}
% The use of low-rank matrix decomposition for compressing Deep Neural Networks (DNNs) represents a straightforward yet effective approach, garnering considerable attention within both scientific computing and machine learning domains. In recent years, the challenge of efficiently compressing and accelerating large-scale neural networks via low-rank methods has become a focal point of research. This has led to significant advancements in developing and refining low-rank factorization strategies tailored for DNNs~\citep{schotthofer2022low}.
In the realm of low-rank decomposition~\citep{schotthofer2022low} for neural network compression, existing methods can be broadly classified into two categories: \textbf{fixed low rank} and \textbf{variable low rank} approaches. Fixed rank methods typically involve decomposing weight matrices of pre-trained networks using techniques like Singular Value Decomposition (SVD) or tensor decomposition, followed by fine-tuning the factorized network~\citep{denton2014exploiting,lebedev2014speeding,sainath2013low,moczulski2015acdc}. They also involve constraining weight matrices to maintain a fixed low rank during training~\citep{jaderberg2014speeding,khodak2021initialization,wang2021pufferfish}, or constructing layers as linear combinations of layers with varying ranks~\citep{ioannou2015training}. 
A notable limitation of these methods is the introduction of matrix decomposition rank as a hyperparameter requiring fine-tuning. In contrast, rank-adaptive methods address this limitation by automatically determining and adjusting the low-rank structure. In particular, \citet{kim2015compression,kim2019efficient} apply heuristics search to pre-determine the decomposition rank, while \citet{wen2017coordinating} learn low-rank weights through a loss function penalizing approximated matrix ranks. 
\citet{li2023losparse} use low-rank approximation plus a sparse matrix to compress the weight matrix in transformers. 
% These approaches remove the need for selecting hyper-parameters of LLMs and highlight the potential of low-rank decomposition as an effective tool for model compression, paving the way towards more efficient and scalable LLMs. 

However, none of these methods have worked in the era of LLMs due to their training-require nature. We propose \asvd, a \textit{post-training} LLM decomposition approach enabling the adaptive determination of SVD ranks to optimize the matrix approximations based on feature activations. 
To our knowledge, \asvd~represents the first attempt to compress the weights of LLMs through decomposition in a training-free manner. Since the introduction of \asvd, there have been subsequent works on training-free LLM decomposition, such as SVD-LLM \citep{wang2024svdllm} and Palu~\citep{chang2024palu}. 
% These follow-up studies underscore the significance and potential of our approach. We hope that our proposed post-training LLM decomposition method can establish a new paradigm for LLM compression, opening up avenues for more efficient and accessible deployment of LLMs.